%───────────────────────────────────────────────────────────────────────────────
\subsection{MOTFM Baseline}
\label{sec:motfm_baseline}

Medical Optimal Transport Flow Matching (MOTFM)~\cite{motfm} is our
primary baseline: a conditional flow matching model operating in pixel space
with an ODE solver at inference.  Unlike NeuroMF, MOTFM requires multiple
integration steps (typically 10 Euler steps) and works directly on voxel-space
volumes rather than in a compressed latent space.

\subsubsection{Architecture}

MOTFM uses a \code{DiffusionModelUNet} from MONAI-Generative with 4 resolution
levels: channel widths $[32, 64, 128, 256]$, 2 residual blocks per level,
attention only at the coarsest level, and flash attention enabled.  The model is
unconditional (no mask or class conditioning) for fair comparison with NeuroMF
on the same dataset.

\subsubsection{Step-Matched Training Budget}

To ensure a fair comparison, we match the total number of \emph{optimiser steps}
rather than epochs, since our dataset is $7.3\times$ larger than the one used in
the original MOTFM paper.

\begin{table}[ht]
\centering
\caption{MOTFM training budget: paper vs.\ our step-matched setup.}
\label{tab:motfm_budget}
\begin{tabular}{lcc}
\toprule
\textbf{Property} & \textbf{Paper} & \textbf{Ours} \\
\midrule
Dataset                  & MSD Brain (750 scans) & FOMO-60K (5{,}471 scans) \\
Resolution               & $192^3$               & $192^3$ \\
Hardware                 & 1$\times$ RTX 4090    & 4$\times$ A100-40GB (DDP) \\
Batch size (per GPU)     & 1                     & 1 \\
Gradient accumulation    & 1                     & 1 \\
Effective batch size     & 1                     & 4 \\
Steps per epoch          & 750                   & 1{,}368 \\
Epochs                   & 200                   & 110 \\
\textbf{Total optimiser steps} & \textbf{150{,}000} & \textbf{${\sim}$150{,}480} \\
Learning rate            & $10^{-4}$             & $10^{-4}$ \\
Precision                & fp32                  & bf16-mixed \\
Solver                   & Euler, 10 steps       & Euler, 10 steps \\
\midrule
Est.\ wall time          & ---                   & ${\sim}$69\,h (2.9 days) \\
SLURM wall limit         & ---                   & 4 days \\
\bottomrule
\end{tabular}
\end{table}

The step-matching rationale is straightforward:
\begin{equation}
  N_\text{epochs}^\text{ours}
  = \left\lceil \frac{N_\text{paper} \cdot E_\text{paper}}
                     {\lceil S_\text{ours} / G \rceil} \right\rceil
  = \left\lceil \frac{750 \times 200}{\lceil 5471 / 4 \rceil} \right\rceil
  = \left\lceil \frac{150{,}000}{1{,}368} \right\rceil
  = 110
\end{equation}
where $S_\text{ours} = 5{,}471$ is our training set size and $G = 4$ is the GPU
count.  The effective batch size increases from~1 to~4 due to DDP data
parallelism; this is standard practice and commonly accepted for fair
comparison, as each GPU still processes one sample per step with no gradient
accumulation.

\subsubsection{Key Differences from NeuroMF}

\begin{table}[ht]
\centering
\caption{Architectural and methodological differences between NeuroMF and MOTFM.}
\label{tab:neuromf_vs_motfm}
\begin{tabular}{lcc}
\toprule
\textbf{Property} & \textbf{NeuroMF} & \textbf{MOTFM} \\
\midrule
Operating space      & Latent ($4 \times 48^3$) & Pixel ($1 \times 192^3$) \\
NFE (inference)      & 1                        & 10 (Euler) \\
Training paradigm    & MeanFlow (iMF dual-head) & Flow matching (velocity) \\
VAE dependency       & Frozen MAISI VAE         & None \\
Model parameters     & ${\sim}$178M             & TBD \\
Conditioning         & Unconditional            & Unconditional \\
\bottomrule
\end{tabular}
\end{table}
